%kódy na zkopírování pro vložení obrázků


%clean verze:
\begin{figure}[htbp]
	\centering
	\includegraphics[width=0.8\textwidth]{nazev_obrazku}
	\caption{Popisek obrázku, který se objeví pod ním a v seznamu obrázků.} 
	\label{fig:muj_obrazek}
\end{figure}



\begin{figure}[htbp]      % [htbp] říká: "Zkus to vložit Tady, nebo Nahoře, Dole, nebo na samostatné Stránce."
    \centering             % Obrázek bude vycentrován na střed stránky
    \includegraphics[width=0.8\textwidth]{nazev_obrazku} % Nastaví šířku na 80 % šířky textu
    \caption{Popisek obrázku, který se objeví pod ním a v seznamu obrázků.} 
    \label{fig:muj_obrazek} % Značka pro křížové odkazy v textu
\end{figure}

%kdyžtak před width jde ještě dát "angle=90, " na řádek 5 pro otočení
%když je obrázek ve složce dám před jeho jméno \jméno_složky


% nebo víc vedle sebe

\begin{figure}[htbp]
    \centering
    \begin{subfigure}[b]{0.45\textwidth} % 45 % šířky textu pro první pod-obrázek
        \centering
        \includegraphics[width=\textwidth]{z_beta}
        \caption{Záření $\beta$ ($^{90}\text{Sr}$).}
    \end{subfigure}
    \hfill % Vloží pružnou mezeru mezi pod-obrázky
    \begin{subfigure}[b]{0.45\textwidth} % 45 % šířky textu pro druhý pod-obrázek
        \centering
        \includegraphics[width=\textwidth]{z_gama}
        \caption{Záření $\gamma$ ($^{137}\text{Cs}$).}
    \end{subfigure}
    \caption{Porovnání útlumu různých typů ionizujícího záření v hliníku.}
    \label{fig:porovnani}
\end{figure}


%sample equation

\begin{equation}
	nR = \frac{p_1 V_1}{T_1} = \frac{p_2 V_2}{T_2}
\end{equation}
\vspace{0.5em}
\noindent kde je: \\
\begin{tabular}{l l l}
	$n$ & látkové množství & $[\text{mol}]$ \\
	$R$ & plynová konstanta & $[\text{J}\cdot\text{K}^{-1}\cdot\text{mol}^{-1}]$ \\
	$p$ & tlak plynu & $[\text{Pa}]$ \\
	$V$ & objem plynu & $[\text{m}^3]$ \\
	$T$ & termodynamická teplota & $[\text{K}]$ \\
\end{tabular}
\vspace{0.5em} 